\chapter{引言}\label{chap:introduction}

\section{研究背景}

NP 完全性将大量组合问题置于多项式时间归约所定义的统一框架中，是比较问题计算难度的重要理论工具。Cook 和 Karp 建立了 NP 完全性与多项式时间归约的基本理论\citep{cook1971complexity,karp1972reducibility}。此后，许多组合问题通过归约获得困难性结论\citep{garey1979computers}。证明一个问题为 NP 完全问题，既要说明该问题属于 NP，也要证明其 NP 困难性。后者通常需要选择一个已知 NP 困难问题，并构造保持判定结果的实例映射。

\begin{figure}[!htbp]
    \centering
    \includegraphics[width=0.88\textwidth]{npc}
    \bicaption{NP 完全性的证明条件。}{Proof requirements for NP-completeness.}
    \label{fig:npc}
\end{figure}

设 \(\Sigma\) 为有限字母表，\(\Sigma^*\) 表示由 \(\Sigma\) 中符号组成的所有有限字符串的集合。复杂性理论用语言 \(L\subseteq\Sigma^*\) 表示判定问题。具体问题的输入对象经过编码后成为字符串，其中属于 \(L\) 的字符串称为 YES 实例，不属于 \(L\) 的字符串称为 NO 实例。将源问题和目标问题对应的语言分别记为 \(L_1,L_2\subseteq\Sigma^*\)。若存在多项式时间可计算函数 \(f:\Sigma^*\rightarrow\Sigma^*\)，满足

\begin{equation}
    \forall x\in\Sigma^*,\qquad
    x\in L_1 \Longleftrightarrow f(x)\in L_2,
    \label{eq:many-one-reduction}
\end{equation}

则称 \(L_1\) 可以多项式时间多对一归约到 \(L_2\)，记作 \(L_1\leq_m^p L_2\)。当 \(L_1\) 对应的问题已知为 NP 困难问题时，式~\eqref{eq:many-one-reduction}可用于证明 \(L_2\) 对应问题的 NP 困难性。完整证明需要给出映射 \(f\) 的构造，并分别证明 \(x\in L_1\Rightarrow f(x)\in L_2\) 和 \(f(x)\in L_2\Rightarrow x\in L_1\)。运行时间与输出长度还须受输入长度的多项式限制。

\begin{figure}[!htbp]
    \centering
    \includegraphics[width=0.88\textwidth]{reduction}
    \bicaption{多项式时间归约关系。}{Basic relation in polynomial-time reduction.}
    \label{fig:reduction}
\end{figure}

归约构造的难点不在于复述定义，而在于寻找满足全部条件的映射。源问题选择失当参数变换错误，都可能使证明失败。候选空间通常难以穷举，错误又常在边界实例分析或证明展开时才会暴露。研究者需要反复检索已有结论、比较源问题、实现候选映射、分析实例并补全证明，人工经验难以系统覆盖这一搜索过程。

人工智能在数学推理中的应用正由答案生成转向可核验的证明构造，并已用于形式化证明、欧氏几何证明和数学构造问题\citep{zheng2022minif2f,yang2023leandojo,trinh2024alphageometry,romeraparedes2024funsearch,hubert2026alphaproof}。这些工作表明，人工智能可以生成数学对象并依据外部反馈加以检验。但归约发现不仅需要证明给定命题，还要选择源问题并设计新的实例映射。

大语言模型可以依据问题定义和已有知识生成候选映射，为归约搜索提供起点。由于单次生成难以完成检索、构造、检验和证明修正的全过程，本课题同时结合大语言模型和智能体环境，使其依据各阶段反馈持续修正归约证明，最终以期得到正确的归约。

\section{研究意义}

本课题的首要意义在于为复杂度尚未确定的 NP 判定问题寻找新的复杂性证明。若能构造满足全部条件的多项式时间归约，便可证明目标问题的 NP 困难性；结合该问题属于 NP 的证明，还可得到新的 NP 完全性结论。这类结果能够扩充已知 NP 完全问题及其归约关系，并揭示目标问题与经典困难问题之间的结构联系。无法得到完整归约时，受限问题的困难性结论、已经证明的局部构造或能够排除某类候选映射的反例，也可以缩小后续研究的范围。

本课题也探索人工智能辅助理论证明的研究方法。现有人工智能定理证明工作多在给定命题后搜索证明\citep{lample2022hypertree,jiang2023draft,xin2024deepseekprover}，复杂性归约还需要选择源问题并设计新的实例映射。本课题以智能体组织大语言模型，使其依据反馈迭代构造候选，并通过数学证明核验结果，为人工智能参与开放式理论探索提供参考。

已知归约任务还为考察人工智能的推理能力提供了可核对的实验对象。通过改变问题结构、组合归约长度和反馈方式，可以分析模型能否理解问题定义、建立结构对应、利用执行与证明反馈并完成多步映射。相关结果用于刻画给定模型、运行环境和计算资源条件下的能力边界，区分可由外部反馈改善的推理环节与仍需依赖新结构发现的环节。
